\begin{abstract}
Version~6 is an artifact-verified revision of the v5 audit of the
constant-transverse-profile calculation for the scoped O.1 problem
on \(C_\alpha\times\RR^2\).
It keeps v5's structure, its evidence classes, and its conditional
theorems. What it changes is that the executable record behind v4
is now archived, rerun, and audited, and it adds the following.

\smallskip\noindent\textbf{New in v6.}
\begin{enumerate}[label=\textup{(\arabic*)},leftmargin=1.6em]
\item \emph{Catalogue re-derived.}
Every entry of the v5 coefficient catalogue agrees exactly with the
executed spectral reductions of the frozen XVI-P vertex arrays.
The executed record also contains three things v4 never printed:
the \((12)(12)\) entries
\(\widehat B_{\mathrm{D01}}=\widehat B_{\mathrm{D03}}
=-8c_1P^2\,\ka^2/(16\pi)\); the complete D02 tadpole array
\(\tfrac{\ka^2}{2\pi}(c_0+2c_1+c_2)\) times a fixed tensor;
and \(\widehat B_{\mathrm{D12}}=0\).
Mixed-helicity monomials occur only on vector-type keys, which are
excluded from the record, not set to zero.
The D13 transverse source trace is computed to be zero.

\item \emph{The \(hh\) Slavnov--Taylor failure is diagnosed.}
The v4 failure comes from the coordinate weight of the antifield
\(K'_{(03)}\). The correct weight is \(-2\); v4 scanned only
positive weights.
With conventions derived from the frozen vertex, the
constant-profile \(hh/Kc\) row closes on CR and on F.
The closure is exact as a rational identity in \(\alpha\), with
\(\lambda=-2\ii\). This value is fixed in advance by the tree-level
data, not fitted. UNIT closes trivially.

\item \emph{A second apex coefficient.}
We prove the \(r^2\)-weighted apex coefficient
\[
\mathcal D_b=-\frac{11\alpha}{180}+\frac{B_2(b)}{3\alpha}
-\frac{B_4(b)}{6\alpha^3},
\qquad
\ApTr\bigl[r^2f(A_b)\bigr]=-\mathcal D_bf'(0).
\]
Its singular-channel shift is \(2\nu(1-\nu^2)/3\).
It passes exact orbifold and twisted-orbifold checks, and it has no
flat CR remnant.

\item \emph{An exact obstruction for \(\HI\) in D01, and its
resolution.}
The \(S=2\to3\) composite trace has the non-integrable diagonal
\(r^{-2\nu-2}\), so \(\HI\) fails as an \(L^2\) statement.
Its unique homogeneous (Mellin) continuation, however, vanishes
identically in \(\nu\). This is shown in closed form and confirmed
independently by quadrature.

\item \emph{Invariant structure of the \(hh\) poles.}
The constant-profile \(hh\) poles decompose into six
\(SO(1,1)\)-invariant structures.
In the complete \(hh\) sum the purely longitudinal structures
cancel, and the gauge variance enters only through
\(b_6+b_7=-\ka^2V/(4\pi)\).
The CR--F difference is gauge invariant:
\(-\nu\ka^2P^2f_0f_0'/(4\pi)\).
\end{enumerate}
A point-by-point table answers every issue v5 raised against v4.
An appendix lists full SHA-256 digests and rerun commands.

\smallskip\noindent\textbf{Still open.}
The 26 Ward combinations and the 34-dimensional projection;
the six exact values;
the D04 quadratic-profile part of the second \(\phi\phi\) value;
the evanescent coefficients;
the 18-dimensional engine map;
an exact \(\Span\) certificate;
and the vector-type keys.
None of these is reported as zero.
The scope remains one loop at \(O(\ka^2)\), the massless two-point
sector, frozen de Donder gauge, and the XVI-P prescription.
The cone and the CR branch remain inputs, and the
destructive-readout analogy supplies no premise.
\(\texttt{bulk\_symmetry\_restored}=\texttt{false}\), and
O.1 on the cone remains \(\NOTCLOSED\).
\end{abstract}
